%=============================================================================
%  07_negative_results.tex -- Negative results
%
%  CONVENTION: as in 06_applications.tex. Measured quantities are macros from
%  macros.tex, or \needsnum{suggestedMacroName} where no macro exists yet.
%  Sweep points and sample sizes are design parameters and appear as literals.
%=============================================================================
\section{Negative Results}
\label{sec:negative}

Five of the hypotheses we set out to test failed, and one of them failed in
every category we measured. We report them here at full length rather than in a
limitations paragraph, for two reasons. First, three of the five are more
informative than the results that passed: a mechanism that fails under one task
framing and succeeds under another (\cref{sec:neg-critic}) tells us more about
masked diffusion than either observation alone. Second, the strongest single
result in this paper --- the quality invariance of \cref{sec:cost-laws} --- was
found while investigating a claim that turned out to be false, and we would not
have found it if we had stopped reporting at the first refutation.

%-----------------------------------------------------------------------------
\subsection{Perturbation robustness: refuted in every category}
\label{sec:neg-perturb}

\paragraph{Prediction.} Masked diffusion denoises the whole output canvas
jointly, conditioned bidirectionally on the input. We predicted this would make
it more robust to \emph{input} noise than left-to-right decoding: a corrupted
token in the log line would be outvoted by the surrounding context during
denoising, whereas AR, having committed to a prefix, would carry the corruption
forward. Log streams drift --- formats change, fields get abbreviated,
separators change with a library upgrade --- so robustness to input perturbation
is the property that would matter most in production.

\paragraph{Measurement.} We apply \perturbCategories{} perturbation categories
at intensity $1.0$ to the input log lines and measure Field-F1 against the
unperturbed gold records. Clean baselines are \perturbCleanAR{} for AR and
\perturbCleanDM{} for diffusion; the drops in \cref{tab:perturb} are relative to
each model's own clean baseline, so the comparison is of \emph{degradation},
not of level.

\begin{table}[t]
  \centering
  \small
  \caption{Perturbation robustness at intensity $1.0$, as F1 drop from each
  model's own clean baseline (AR \perturbCleanAR, diffusion \perturbCleanDM).
  Diffusion degrades more in \perturbCategories{} of \perturbCategories{}
  categories. Rows are ordered by the ratio of drops.}
  \label{tab:perturb}
  \begin{tabular}{lccc}
    \toprule
    Perturbation & AR drop & diffusion drop & ratio\\
    \midrule
    whitespace            & \needsnum{pertDropARwhitespace}  & \needsnum{pertDropDMwhitespace}  & $\infty$ (AR flat)\\
    abbreviation          & \needsnum{pertDropARabbrev}      & \needsnum{pertDropDMabbrev}      & $\infty$ (AR flat)\\
    separator\_change     & \needsnum{pertDropARsep}         & \needsnum{pertDropDMsep}         & \perturbWorstRatio$\times$\\
    encoding\_noise       & \needsnum{pertDropARencoding}    & \needsnum{pertDropDMencoding}    & \needsnum{pertRatioEncoding}$\times$\\
    wrapping\_noise       & \needsnum{pertDropARwrapping}    & \needsnum{pertDropDMwrapping}    & \needsnum{pertRatioWrapping}$\times$\\
    missing\_field        & \needsnum{pertDropARmissing}     & \needsnum{pertDropDMmissing}     & \needsnum{pertRatioMissing}$\times$\\
    field\_reorder        & \needsnum{pertDropARreorder}     & \needsnum{pertDropDMreorder}     & \needsnum{pertRatioReorder}$\times$\\
    timestamp\_format     & \needsnum{pertDropARtimestamp}   & \needsnum{pertDropDMtimestamp}   & \needsnum{pertRatioTimestamp}$\times$\\
    synonym\_substitution & \needsnum{pertDropARsynonym}     & \needsnum{pertDropDMsynonym}     & \needsnum{pertRatioSynonym}$\times$\\
    case\_variation       & \needsnum{pertDropARcase}        & \needsnum{pertDropDMcase}        & \needsnum{pertRatioCase}$\times$\\
    \bottomrule
  \end{tabular}
\end{table}

\paragraph{Result.} The prediction is refuted in \perturbCategories{} of
\perturbCategories{} categories. Diffusion degrades more than AR under every
perturbation we applied. The worst case is \texttt{separator\_change}, where
diffusion's drop is \perturbWorstRatio$\times$ AR's. On two categories ---
\texttt{whitespace} and \texttt{abbreviation} --- AR is \emph{flat}, its drop
indistinguishable from zero, while diffusion loses measurable F1 on both. The
result is not marginal, not confined to a subset of perturbation types, and not
in the direction we predicted.

\paragraph{Hypothesis, offered as a hypothesis.} One explanation fits the shape
of the data: in left-to-right decoding, a corrupted input token influences every
output position that follows it and none that precede it, so the damage is
bounded by the position at which the corruption enters. Under bidirectional
conditioning every output position attends to every input position at every
denoising step, so a single corrupted token can influence the entire canvas, and
each denoising iteration re-exposes the whole record to it. Noise propagates
forward only under causal attention; it propagates in both directions under
bidirectional attention.

We emphasise that this is a hypothesis and not a conclusion. It is not the only
one available: the diffusion model here is \modelLLaDA{}, whose pretraining
corpus and scale differ from \modelARthree{}'s, and a pretraining artifact would
produce the same table. Distinguishing them requires a size-matched ablation ---
the same perturbation suite run against \modelDG{} --- which we have not run.
The mechanism is unresolved and tracked as issue~\#39. \open{\#39}

%-----------------------------------------------------------------------------
\subsection{One-shot imputation: refuted}
\label{sec:neg-impute}

\paragraph{Prediction.} Filling in $K$ missing fields of a partially observed
record is masked diffusion's defining operation: place mask tokens at the
missing field positions, denoise, read off the result. AR has no architectural
answer --- it must regenerate the record and be trusted to copy the observed
fields back verbatim. We predicted diffusion would win this task outright, and
that its cost would be flat in $K$ while AR's grew.

\paragraph{Measurement.} We mask $K \in \{1,2,4,6,8\}$ fields of a
twelve-field record and evaluate F1 on the masked fields only.

\begin{table}[t]
  \centering
  \small
  \caption{One-shot imputation of $K$ masked fields, F1 on the masked fields.
  Diffusion loses at every $K$, and by the widest margin where its mechanism
  should help most.}
  \label{tab:impute}
  \begin{tabular}{cccc}
    \toprule
    $K$ masked & AR F1 & diffusion F1 & verdict\\
    \midrule
    1 & \imputeARLo                    & \imputeDMLo                    & AR\\
    2 & \needsnum{imputeARKtwo}        & \imputeDMHi                    & AR\\
    4 & \imputeARHi                    & \needsnum{imputeDMKfour}       & AR\\
    6 & \needsnum{imputeARKsix}        & \needsnum{imputeDMKsix}        & AR\\
    8 & \needsnum{imputeARKeight}      & \needsnum{imputeDMKeight}      & AR\\
    \bottomrule
  \end{tabular}
\end{table}

\paragraph{Result.} Diffusion scores \imputeDMLo--\imputeDMHi{} across the
sweep against AR's \imputeARLo--\imputeARHi. At $K=1$ --- a single masked field,
the cleanest possible statement of the mechanism --- diffusion returns
\imputeDMLo. The cost predictions fare no better: diffusion's latency is indeed
flat in $K$ (ratio \needsnum{imputeLatencyRatio} across the sweep), but so is
\emph{AR's}, at the same ratio, because the record length does not change with
$K$; and the cost ratio between the two models is flat at
\needsnum{imputeCostRatio}$\times$ for all $K$ rather than widening. Of the four
sub-predictions, one passes and three are refuted, and the one that passes
(diffusion latency flat in $K$) is not discriminative because AR shares it.

%-----------------------------------------------------------------------------
\subsubsection{Task framing beat mechanism: the critic is the differentiator}
\label{sec:neg-critic}

The two masking experiments are the same code path. \Cref{sec:neg-impute} masks
$K$ fields and asks the model to fill them given only the log line; iterative
refinement (\cref{sec:app-refine}) masks the fields a critic has flagged and
asks the model to fix them given the rest of the record. The mask-construction
routine, the denoising loop, and the model weights are identical. The first
returns \imputeDMLo--\imputeDMHi{} F1; the second gains \refineGainDM{} F1
over three rounds.

The difference is not the mechanism. It is what selects the mask and what
context survives it. Under refinement the model conditions on the fields the
critic left alone --- a partially correct record is a much stronger conditioning
signal than a bare log line --- and the critic's selection concentrates the
model's effort on positions where it is likely to be wrong. Under one-shot
imputation the selection is uninformed and the surviving context is weak.

We take this seriously as a finding about how to \emph{use} masked diffusion:
the native mask operation is not by itself a capability, and framing the task so
that the mask lands where the model can act on it is what produces the gain.
It is not yet a confirmed causal claim. The ablation that would confirm it is
straightforward and unrun: hold the refinement loop fixed and replace the critic
with a random selection of the same number of fields. If the gain survives, the
conditioning context is doing the work; if it disappears, the critic's selection
is. We propose that ablation rather than assert its outcome.

%-----------------------------------------------------------------------------
\subsection{Adaptive effort: mixed}
\label{sec:neg-adaptive}

\paragraph{Prediction.} Adaptive threshold decoding stops denoising when
confidence is high, so a diffusion model should spend fewer steps on routine
events and more on novel ones --- fast where the input is easy, deep where it is
hard --- while AR pays the same per-token cost regardless.

\paragraph{Measurement.} We bin BGL events into routine, medium, and novel by
template frequency and measure both models per bin. The quality half of the
prediction holds, and holds in an interesting direction: diffusion reaches
\adaptDMNovel{} F1 on novel events against AR's \adaptARNovel, and diffusion's
F1 \emph{improves} from routine (\needsnum{adaptDMRoutine}) to novel
(\adaptDMNovel) while AR's degrades over the same gradient
(\needsnum{adaptARRoutine} to \adaptARNovel). A model that gets better on the
rare tail is unusual and worth following up.

\paragraph{The latency half is refuted, and it is not close.} Diffusion's $p50$
is \adaptDMPFifty~ms against AR's \adaptARPFifty~ms --- roughly
\needsnum{adaptSlowdown}$\times$ slower. The adaptive decoder buys its quality
with an order of magnitude of latency, and does so on a runtime whose KV-cache
is broken for our prompt lengths (\cref{sec:threats}), so this latency figure is
a property of that configuration rather than of adaptive decoding as such. The
``fast on easy inputs'' half of the prediction did not materialise at all.

\paragraph{Caveat.} The difficulty bins use exact template-match frequency as a
proxy for difficulty. Events binned as novel may be novel only in the sense of
having no exact template sibling in this corpus, which is not the same as being
semantically harder; a semantic clustering of templates would be needed to rule
out a binning artifact, and that check is open. Sample sizes also differ between
arms (100 per bin for diffusion against 200 per bin for AR).

%-----------------------------------------------------------------------------
\subsection{Semi-autoregressive blocks: a schema mechanism, not a general trick}
\label{sec:neg-semiar}

\paragraph{Prediction.} Decoding the canvas in $B$ sequential blocks rather than
all at once should restore some of the left-to-right conditioning that
cross-field consistency depends on, recovering consistency without giving up
diffusion's cost profile.

\paragraph{Measurement.} Consistency does recover, sharply: from
\semiARConsistBOne\% at $B=1$ to \semiARConsistBTwo\% at $B=2$. Field-F1
collapses over the same change, from \semiARFOneBOne{} to \semiARFOneBTwo. At
$B=4$ both metrics are worse than at $B=2$ (\needsnum{semiARFOneBFour} F1,
\needsnum{semiARConsistBFour}\% consistency). Latency is essentially unchanged
across the sweep at \needsnum{semiARPFifty}~ms.

\paragraph{Interpretation.} Our block boundaries fall at raw token positions,
not at field boundaries, so a block split routinely cuts a field in half. A
model trained with explicit schema slots would place its block boundaries where
the fields are; an untuned model has no slot priors and cannot. What the sweep
measures, then, is that consistency improves whenever the decoder is forced to
commit to part of the canvas before the rest --- and that doing so at the wrong
positions destroys the content. Semi-autoregressive blocking is a
schema-template mechanism, inseparable from schema-aware training, and not a
general trick that can be applied to an off-the-shelf masked diffusion model.
This ablation is preliminary: $n=30$.

%-----------------------------------------------------------------------------
\subsection{The classical-baseline ceiling}
\label{sec:neg-baselines}

The most consequential negative result in this paper is not about diffusion at
all. It is that a 2017 log parser does most of this task better than either
language model, thousands of times faster.

On BGL, Drain alone reaches \drainFOne{} Field-F1 and \drainAnomalyFOne{}
anomaly-F1 at \drainEvs{} events per second. Per field, it scores
\drainStructuralLo--\drainStructuralHi{} on \emph{all seven} structural fields
--- timestamp, host, component, severity, template, error code, metric value ---
matching or exceeding both language models, which reach
\needsnum{llmStructuralPerField} on the same fields. A rule-based extractor
reaches \ruleFOne{} overall, well behind Drain but at comparable speed.
One implementation caveat belongs here rather than in a footnote: the reference
\texttt{logparser} package was unavailable in the measurement environment, so
our Drain numbers come from a regex-fallback implementation of the same
fixed-depth-tree algorithm. If anything this understates the baseline, which
makes the conclusion below stronger rather than weaker --- but the numbers
should be reproduced against the reference implementation before publication.

The consequence for how this task should be read is direct. Of the twelve
fields in our schema, seven are recoverable by template mining at or above LLM
accuracy. The language model's contribution is confined to the three semantic
fields, which no parser can produce, and the two anomaly fields, on which Drain
already matches the LLMs when native labels exist. Any end-to-end structured
extraction result reported without this baseline attributes to the model a
majority of the score that a parser would have produced for free, and therefore
oversells. We state that plainly because our own earlier framing of this work
did exactly that.

There is a real boundary, and it is the one negative result in this subsection
that runs the other way. Under mixed perturbation at intensity $1.0$, Drain
degrades to \drainPerturbed{} and the rule-based extractor to \rulePerturbed,
while the language models hold much closer to their clean baselines
(\cref{tab:perturb}). Template mining depends on the templates being stable;
when the input format drifts, the classical methods lose what makes them work
and the language models do not. That --- adversarial or drifting log
environments --- is where an LLM earns its place, and it is a much narrower
claim than ``LLMs parse logs well.''

The scale argument is equally blunt. On Thunderbird's 211M-line corpus, the
diffusion pipeline would take an estimated \thunderbirdDM{} years; Drain
finishes in \thunderbirdDrain{} minutes. There is no serving optimisation that
closes a gap of that size. The realistic architecture for log extraction is a
parser carrying the structural fields with a language model invoked selectively
on the semantic ones, and the interesting engineering question is when to
invoke it --- not which architecture to invoke.

%-----------------------------------------------------------------------------
\subsection{What survived: diffusion-internal claims}
\label{sec:neg-survived}

Not everything failed, and the pattern in what did not is worth stating
explicitly. Latency is linear in the number of denoising steps $K$ with
$R^2 = \kSweepRsq$ and a slope of \kSweepSlope~ms per step; the quality knee
sits at $K = \kKneeLo$--$\kKneeHi$, with F1 rising steeply from $K=2$
(\needsnum{kSweepFOneKtwo}) to $K=8$ (\needsnum{kSweepFOneKeight}) and barely
moving to $K=24$ (\needsnum{kSweepFOneKtwentyfour}); and cross-field consistency
is the model's identified weak spot at \mdmConsistency\% against AR's
\arConsistency\%.

Every one of these is a \emph{diffusion-internal} claim: a statement about how
one model behaves as one of its own parameters varies, verifiable without
reference to any other architecture. Every claim we refuted, in this section and
in \cref{sec:applications}, is \emph{comparative}: a statement that diffusion
beats AR at something. The internal claims survived contact with a corrected
baseline; the comparative ones largely did not. That distinction is the
methodological result of this paper, and \cref{sec:related} places it against
the reproducibility literature that predicts exactly this failure mode.
